\otazka{5. Kyberkriminalita}{Prameny práva (národní, evropské a mezinárodní), typická trestná činnost, klasifikace trestných činů, právní kvalifikace a související postupy a kritéria, mezinárodní spolupráce.}{BVV03K}

\subsection*{5. Kyberkriminalita -- stručněji}
Kyberkriminalita je trestná činnost, kde jsou technologie cílem, nástrojem nebo prostředím. U zkoušky je dobré jít přes prameny práva, typické skutky, klasifikaci, paragrafy, procesní postupy a mezinárodní spolupráci.

\textbf{1. Prameny práva}
\begin{itemize}
  \item Národní rámec tvoří hlavně trestní zákoník a trestní řád. Trestní zákoník říká, co je trestné, trestní řád říká, jak se trestná činnost vyšetřuje a dokazují se elektronické stopy.
  \item Evropsky jsou důležité předpisy k útokům na informační systémy, platebním podvodům, evropskému vyšetřovacímu příkazu a e-evidence. Mezinárodně je základ Budapešťská úmluva.
\end{itemize}

\textbf{2. Typická trestná činnost}
\begin{itemize}
  \item Cyber-dependent kriminalita existuje jen díky IT: neoprávněný přístup, malware, DDoS, ransomware, zásah do systému nebo dat.
  \item Cyber-enabled kriminalita je tradiční trestná činnost usnadněná internetem: podvod, vydírání, krádež identity, phishing, platební podvody, stalking nebo šíření nelegálního obsahu.
  \item Cyber-supported činnost používá IT jen podpůrně, například komunikace pachatelů, ukládání výnosů nebo koordinace skupiny.
\end{itemize}

\textbf{3. Právní kvalifikace}
\begin{itemize}
  \item Typické paragrafy jsou § 230 neoprávněný přístup k počítačovému systému a nosiči informací, § 231 opatření a přechovávání přístupového zařízení a hesla, § 232 poškození záznamu v počítačovém systému.
  \item Podle situace se přidává § 209 podvod, § 175 vydírání, § 234 neoprávněné opatření platebního prostředku nebo § 182 porušení tajemství dopravovaných zpráv.
  \item Kvalifikace závisí na jednání, úmyslu, škodě, rozsahu, následku, souběhu trestných činů a tom, zda pachatel útočil na systém, data, peníze nebo člověka.
\end{itemize}

\textbf{4. Související postupy}
\begin{itemize}
  \item U kyberkriminality je klíčové rychle zajistit elektronické důkazy, protože logy mizí, data se mění a služby jsou často v cloudu.
  \item Prakticky se řeší uchování dat, vydání nebo odnětí zařízení, metadata, domovní prohlídka, forenzní kopie, hash, chain of custody a spolupráce s poskytovateli.
  \item Podle zákona o elektronických komunikacích operátoři veřejných sítí nebo veřejně dostupných služeb uchovávají provozní a lokalizační údaje 6 měsíců. Jde o metadata, ne obsah zpráv.
\end{itemize}

\textbf{5. Mezinárodní spolupráce}
\begin{itemize}
  \item Kyberkriminalita je často přeshraniční: pachatel, oběť, server, doména, platební tok a poskytovatel mohou být v různých státech.
  \item Používá se MLA, evropský vyšetřovací příkaz, Budapešťská síť 24/7, Europol/EC3, Eurojust a od 18. 8. 2026 e-evidence nástroje EPOC a EPOC-PR.
\end{itemize}

\subsection*{Kostra odpovědi}
\begin{enumerate}
  \item Definovat kyberkriminalitu.
  \item Projít prameny práva: národní, evropské, mezinárodní.
  \item Vyjmenovat typickou trestnou činnost.
  \item Vysvětlit klasifikaci: cyber-dependent, cyber-enabled, cyber-supported a právní třídění podle skutkových podstat.
  \item Ukázat právní kvalifikaci na konkrétních příkladech.
  \item Popsat související postupy a kritéria dokazování.
  \item Uzavřít mezinárodní spoluprací.
\end{enumerate}

\subsection*{Vymezení}
Kyberkriminalita je trestná činnost, u které je informační technologie buď \textbf{cílem}, \textbf{nástrojem}, nebo \textbf{prostředím}. Cílem je například databáze, server, účet, e-mailová schránka nebo platební systém. Nástrojem je technologie tehdy, když pachatel přes phishing, malware, falešný e-shop nebo deepfake páchá klasickou kriminalitu. Prostředím je internet, cloud, sociální sítě, kryptoměnové služby nebo anonymizační infrastruktura.

Odlišuj tři roviny. \textbf{Kyberbezpečnost} řeší ochranu systémů a služeb, typicky NÚKIB, správci a CSIRT/CERT. \textbf{Kyberobrana} je vojenská a státně-obranná rovina. \textbf{Kyberkriminalita} je trestní odpovědnost pachatele: policie, státní zastupitelství, soud, znalci a mezinárodní spolupráce.

\priklad{Ransomware v nemocnici může být současně trestný čin pachatele, kybernetický bezpečnostní incident nemocnice podle ZKB, porušení zabezpečení osobních údajů podle GDPR a krizový problém zdravotnictví. U zkoušky proto vždy řekni, který režim právě řešíš.}

\subsection*{Prameny práva}
Jdi přesně v pořadí: národní právo, evropské právo, mezinárodní právo.

\textbf{Národní prameny:} základ je \textbf{trestní zákoník} a \textbf{trestní řád}. Trestní zákoník říká, co je trestné. Trestní řád říká, jak se to vyšetřuje a dokazují elektronické stopy.
\begin{itemize}
  \item \textbf{§ 230 TZ} - neoprávněný přístup k počítačovému systému a nosiči informací; odst. 1 typicky průnik po překonání zabezpečení, odst. 2 užití, změna, smazání nebo poškození dat.
  \item \textbf{§ 231 TZ} - opatření a přechovávání přístupového zařízení a hesla; typicky prodej hesel, exploitů nebo nástrojů s úmyslem spáchat útok podle § 230.
  \item \textbf{§ 232 TZ} - poškození záznamu v počítačovém systému z nedbalosti; použitelné jen při zákonném následku a nedbalosti, ne na každý omyl v IT.
  \item \textbf{§ 209 TZ} - podvod; phishing, falešný e-shop, investiční podvod, BEC.
  \item \textbf{§ 175 TZ} - vydírání; ransomware, sextortion, hrozba zveřejněním dat.
  \item \textbf{§ 234 TZ} - neoprávněné opatření, padělání a pozměnění platebního prostředku; platební karty, internetbanking, platební údaje.
  \item \textbf{§ 182 TZ} - porušení tajemství dopravovaných zpráv; neoprávněné zachycení komunikace.
\end{itemize}

\textbf{Evropské prameny:}
\begin{itemize}
  \item \textbf{směrnice o útocích na informační systémy} - harmonizuje nezákonný přístup, zásah do systému, zásah do dat, zachycování a nástroje k útokům,
  \item \textbf{směrnice o bezhotovostních platebních podvodech} - důležitá pro karty, elektronické peníze a internetbanking,
  \item \textbf{evropský vyšetřovací příkaz} - orgán jednoho státu EU žádá orgán jiného členského státu o vyšetřovací úkon,
  \item \textbf{e-evidence} - od 18. 8. 2026 EPOC a EPOC-PR pro vydání nebo uchování uložených elektronických důkazů u poskytovatelů v EU.
\end{itemize}

\textbf{Mezinárodní prameny:} základ je \textbf{Budapešťská úmluva o kyberkriminalitě}. Je prakticky nejdůležitější, protože spojuje hmotněprávní skutky, procesní nástroje a spolupráci včetně sítě 24/7. Druhý dodatkový protokol má posílit přímější přeshraniční spolupráci, ale k 21. 6. 2026 ještě není v platnosti. Úmluva OSN proti kyberkriminalitě také ještě není v platnosti; podpis neznamená ratifikaci.

\subsection*{Typická trestná činnost}
\begin{itemize}
  \item \textbf{neoprávněný přístup} - prolomení účtu, zneužití zranitelnosti, průnik do serveru,
  \item \textbf{malware a ransomware} - zašifrování dat, krádež dat, vydírání, dvojí vydírání při hrozbě zveřejněním,
  \item \textbf{phishing a BEC} - vylákání hesel, platebních údajů nebo převodu peněz přes falešnou autoritu,
  \item \textbf{DDoS} - zahlcení služby, často jako vydírání nebo politicky motivovaný útok,
  \item \textbf{krádež a prodej dat} - databáze zákazníků, přístupové údaje, cookies, tokeny, obchodní tajemství,
  \item \textbf{infostealery a initial access brokering} - krádež přístupů a jejich prodej dalším pachatelům,
  \item \textbf{platební a investiční podvody} - karty, internetbanking, kryptoměny, falešné platformy,
  \item \textbf{online vydírání, stalking a sextortion} - klasické útoky zesílené digitálními kanály,
  \item \textbf{nelegální tržiště a praní výnosů} - prodej malwaru, přístupů, drog, dokladů a kryptoměnové praní.
\end{itemize}

\subsection*{Klasifikace trestných činů}
Nejpraktičtější je klasifikovat podle role technologie:
\begin{longtable}{p{0.23\linewidth}p{0.35\linewidth}p{0.34\linewidth}}
\toprule
\textbf{Typ} & \textbf{Co to znamená} & \textbf{Příklad}\\
\midrule
\textbf{Cyber-dependent} & Bez počítače, sítě nebo dat by trestný čin nedával smysl. & Neoprávněný přístup, malware, DDoS, smazání databáze.\\
\addlinespace
\textbf{Cyber-enabled} & Klasická kriminalita, kterou internet usnadní, zrychlí nebo zvětší. & Phishingový podvod, ransomware jako vydírání, falešné investice, sextortion.\\
\addlinespace
\textbf{Cyber-supported} & Technologie je jen podpůrný prostředek. & Komunikace pachatelů přes šifrovanou aplikaci, cloudové úložiště dokumentů, vyhledávání obětí.\\
\bottomrule
\end{longtable}

Vedle toho vždy existuje \textbf{právní klasifikace}: konkrétní skutková podstata v trestním zákoníku. Jeden technický útok může právně spadat pod více trestných činů. Ransomware je technicky malware, ale právně může být současně § 230, § 175 vydírání, případně § 209 podvod nebo další skutky.

\subsection*{Právní kvalifikace}
Právní kvalifikace znamená převést skutkový děj na konkrétní paragraf. Nestačí říct "hacking" nebo "phishing". Postup:
\begin{itemize}
  \item co pachatel udělal a vůči jakému systému, účtu, datům nebo penězům,
  \item jestli překonal zabezpečení, data jen získal, nebo je změnil/smazal/zašifroval,
  \item jaký měl úmysl: průnik, krádež, vydírání, podvod, prodej přístupů,
  \item jaký vznikl následek: škoda, výpadek, únik dat, zásah do komunikace, počet poškozených,
  \item zda jde o souběh více skutkových podstat.
\end{itemize}

\begin{longtable}{p{0.24\linewidth}p{0.37\linewidth}p{0.31\linewidth}}
\toprule
\textbf{Skutek} & \textbf{Typická kvalifikace} & \textbf{Co zdůraznit}\\
\midrule
\textbf{Průnik do účtu/serveru} & § 230 odst. 1, případně § 230 odst. 2. & U odst. 1 se řeší překonání bezpečnostního opatření; u odst. 2 zásah do dat nebo jejich neoprávněné užití.\\
\addlinespace
\textbf{Ransomware} & § 230 odst. 1 a 2, § 175 vydírání, někdy § 209 podvod nebo § 182. & Typický souběh: průnik, zašifrování, exfiltrace, výhrůžka a výkupné.\\
\addlinespace
\textbf{Phishing na banku} & § 209 podvod, § 234 platební prostředek, případně § 230 při přihlášení do účtu. & Dokaž omyl oběti, dispozici s penězi, škodu a použití přístupů.\\
\addlinespace
\textbf{Prodej hesel/exploitu} & § 231, pokud je dán úmysl spáchat § 230. & Samotný nástroj není automaticky trestný; rozhoduje účel a kontext.\\
\addlinespace
\textbf{DDoS} & Podle dopadu § 230 odst. 2, případně vydírání. & Jde hlavně o dostupnost služby, rozsah a škodu.\\
\bottomrule
\end{longtable}

\pozor{U § 230 odst. 1 nestačí říct "někdo se dostal k datům". Musí jít o neoprávněný přístup po překonání bezpečnostního opatření. Když pachatel data mění, maže, zašifruje nebo neoprávněně užije, řeší se hlavně § 230 odst. 2.}

\subsection*{Související postupy a kritéria}
U kyberkriminality nestačí znát paragraf. Musíš ukázat, jak se skutek prokáže. Typická osa:
\begin{enumerate}
  \item \textbf{Zjištění skutku} - trestní oznámení, incident report, banka, CSIRT, poskytovatel, zahraniční partner.
  \item \textbf{Rychlé uchování stop} - logy, účty, IP adresy, e-maily, cloud, domény, kryptopeněženky; data rychle mizí.
  \item \textbf{Zajištění zařízení a dat} - vydání/odnětí věci, prohlídky, forenzní kopie, hash, chain of custody.
  \item \textbf{Vyžádání údajů} - metadata, provozní a lokalizační údaje, data od poskytovatelů, případně obsah uložené komunikace.
  \item \textbf{Analýza} - časová osa, malware, účty, peněžní toky, IP adresy, přihlášení, vazba mezi online identitou a osobou.
  \item \textbf{Kvalifikace a důkazní hodnocení} - skutková podstata, škoda, úmysl, souběh, použitelnost důkazu.
\end{enumerate}

\priklad{Podle zákona o elektronických komunikacích musí právnická nebo fyzická osoba zajišťující veřejnou komunikační síť nebo poskytující veřejně dostupnou službu elektronických komunikací uchovávat \textbf{6 měsíců provozní a lokalizační údaje}. Jde o metadata, například kdo, kdy, odkud a jak komunikoval; ne o obsah zpráv. V trestním řízení se pak k těmto údajům přistupuje jen přes příslušný procesní titul, typicky režim provozních a lokalizačních údajů.}

Hlavní kritéria:
\begin{itemize}
  \item \textbf{zákonnost} - správný procesní titul,
  \item \textbf{proporcionalita} - zásah odpovídá závažnosti,
  \item \textbf{rychlost} - logy rotují, cloudová data se mažou, kryptoměny se přesouvají,
  \item \textbf{integrita} - hash, forenzní kopie, dokumentace,
  \item \textbf{chain of custody} - kdo s důkazem nakládal,
  \item \textbf{jurisdikce} - kde je pachatel, oběť, server, poskytovatel a dopad.
\end{itemize}

\subsection*{Mezinárodní spolupráce}
Kyberkriminalita je skoro vždy přeshraniční. Oběť může být v ČR, server v Irsku, poskytovatel v USA, doména přes anonymního registrátora a peníze v kryptoměně. Proto je nejprve potřeba určit, \textbf{kdo drží data}, \textbf{v jakém státě jsou}, \textbf{o jaký typ dat jde} a \textbf{jak rychle mohou zmizet}.

Používají se:
\begin{itemize}
  \item \textbf{MLA} - klasická vzájemná právní pomoc mezi státy; pomalejší, ale použitelná i mimo EU,
  \item \textbf{EVP/EIO} - evropský vyšetřovací příkaz; v EU orgán žádá orgán jiného členského státu,
  \item \textbf{Budapešťská síť 24/7} - rychlé uchování dat a kontakt mezi státy,
  \item \textbf{Europol/EC3 a Eurojust} - koordinace operací, analytická podpora a koordinace státních zástupců,
  \item \textbf{EPOC/EPOC-PR} - od 18. 8. 2026 přímější příkaz k vydání nebo uchování uložených elektronických důkazů u poskytovatelů v EU,
  \item \textbf{dobrovolná spolupráce poskytovatelů} - jen v mezích práva a pravidel poskytovatele; nenahrazuje soudní nebo jiný povinný titul, pokud je potřeba.
\end{itemize}

\pravobox{EPOC není živý odposlech a neprolamuje šifrování. EPOC slouží k vydání uložených elektronických důkazů, EPOC-PR k jejich uchování. U otázky 5 stačí říct, že jde o evropský nástroj pro rychlejší přeshraniční důkazy u poskytovatelů; detail patří hlavně do otázky 6.}

\zkouska{Souvisle: kyberkriminalita je trestná činnost, kde jsou technologie cílem, nástrojem nebo prostředím. Prameny jsou národní trestní zákoník a trestní řád, evropské směrnice a procesní nástroje a mezinárodně hlavně Budapešťská úmluva. Typicky řeším hacking, malware, ransomware, phishing, DDoS, krádež dat, platební podvody, vydírání a praní výnosů. Klasifikuji cyber-dependent, cyber-enabled a cyber-supported. Potom technický popis převedu na paragrafy: § 230, § 231, § 232, § 209, § 175, § 234, § 182. U kvalifikace řeším jednání, úmysl, škodu, následek a souběh. Procesně je klíčové rychle uchovat elektronické důkazy, zajistit zařízení a data, udržet integritu a řešit jurisdikci. Mezinárodně použiji MLA, EVP, Budapešť 24/7, Europol/Eurojust a nově e-evidence.}

\zdroje{Trestní zákoník; trestní řád; směrnice o útocích na informační systémy; směrnice o bezhotovostních platebních podvodech; směrnice o evropském vyšetřovacím příkazu; Budapešťská úmluva; e-evidence nařízení a směrnice; referenční PDF a starší výpisky.}

\newpage

